\subsection{Generalization Across Model Architectures}
\label{sec:architecture_generalization}

To validate that our linear mergeability prediction framework is not specific to a single model architecture, we replicate all experiments on ViT-B-32 in addition to the primary ViT-B-16 results. Both models share the same Vision Transformer architecture but differ in patch size (16$\times$16 vs 32$\times$32), resulting in different feature dimensions and computational characteristics.

\subsubsection{Validation Performance Comparison}

Table~\ref{tab:arch_comparison} presents the validation Pearson correlation for L1-regularized LOTO across all five merging methods on both architectures.

\begin{table}[h]
\centering
\caption{Validation Pearson correlation ($r$) for L1 LOTO on ViT-B-16 and ViT-B-32. ViT-B-32 achieves higher correlations across all mergers.}
\label{tab:arch_comparison}
\begin{tabular}{l|cc|c}
\toprule
\textbf{Method} & \textbf{ViT-B-16} & \textbf{ViT-B-32} & \textbf{$\Delta$} \\
\midrule
Weight Avg & 0.589 & 0.699 & +0.110 \\
Arithmetic & 0.445 & 0.619 & +0.174 \\
TSV & 0.627 & 0.692 & +0.065 \\
TIES & 0.463 & 0.543 & +0.080 \\
DARE & 0.596 & 0.684 & +0.088 \\
\midrule
\textbf{Average} & 0.544 & \textbf{0.648} & +0.104 \\
\bottomrule
\end{tabular}
\end{table}

ViT-B-32 achieves higher validation correlations across all merging methods, with an average improvement of $+0.104$. This suggests that mergeability prediction is easier for models with larger patch sizes, likely due to simpler feature spaces.

\subsubsection{Metric Selection Analysis}

A critical test of framework generalization is whether the same metrics are identified as important across architectures. Table~\ref{tab:metric_comparison} compares key metric selection statistics.

\begin{table}[h]
\centering
\caption{Comparison of metric selection by L1 LOTO across architectures.}
\label{tab:metric_comparison}
\begin{tabular}{l|cc}
\toprule
\textbf{Statistic} & \textbf{ViT-B-16} & \textbf{ViT-B-32} \\
\midrule
Total metrics & 28 & 28 \\
Universally selected (100\% freq) & 2 & 1 \\
Core predictors (stable + $\geq$80\%) & 2 & 3 \\
Average nonzero per fold & 16--18 & 17--20 \\
\bottomrule
\end{tabular}
\end{table}

\paragraph{Universal Metric: encoder\_gradient\_l2\_distance.}
The most striking finding is that \texttt{encoder\_gradient\_l2\_distance} is the \emph{only} metric universally selected (100\% frequency across all mergers and all folds) on \emph{both} architectures. This metric measures the L2 distance between encoder gradients of two task-specific models, capturing how differently models respond to input perturbations.

Table~\ref{tab:encoder_grad_coef} shows that this metric consistently receives negative coefficients across all mergers on both architectures, indicating that \emph{lower gradient distance predicts better merging outcomes}.

\begin{table}[h]
\centering
\caption{Coefficient for \texttt{encoder\_gradient\_l2\_distance} across mergers. Negative values indicate lower distance predicts better merging. Selection frequency is 100\% for all cells.}
\label{tab:encoder_grad_coef}
\begin{tabular}{l|cc}
\toprule
\textbf{Merger} & \textbf{ViT-B-16 coef} & \textbf{ViT-B-32 coef} \\
\midrule
Weight Avg & $-$0.025 & $-$0.031 \\
Arithmetic & $-$0.016 & $-$0.023 \\
TSV & $-$0.026 & $-$0.028 \\
TIES & $-$0.013 & $-$0.013 \\
DARE & $-$0.026 & $-$0.035 \\
\midrule
\textbf{Average} & $-$0.021 & $-$0.026 \\
\bottomrule
\end{tabular}
\end{table}

\paragraph{Top-1 Metric Agreement.}
For each merger, we identify the metric with the highest importance score (|coefficient| $\times$ selection frequency). Table~\ref{tab:top1_agreement} shows that \texttt{encoder\_gradient\_l2\_distance} is the top predictor for 3 out of 5 mergers on both architectures.

\begin{table}[h]
\centering
\caption{Top-1 metric by importance score per merger. Check marks indicate the metric matches across architectures.}
\label{tab:top1_agreement}
\begin{tabular}{l|ll|c}
\toprule
\textbf{Merger} & \textbf{ViT-B-16 Top-1} & \textbf{ViT-B-32 Top-1} & \textbf{Match} \\
\midrule
Weight Avg & encoder\_grad\_l2 & encoder\_grad\_l2 & \checkmark \\
Arithmetic & input\_grad\_l2 & encoder\_grad\_l2 & \\
TSV & encoder\_grad\_l2 & encoder\_grad\_l2 & \checkmark \\
TIES & interaction\_overlap\_bot & activation\_mag\_ratio & \\
DARE & encoder\_grad\_l2 & encoder\_grad\_l2 & \checkmark \\
\midrule
\multicolumn{3}{r|}{\textbf{Agreement Rate}} & 3/5 (60\%) \\
\bottomrule
\end{tabular}
\end{table}

\paragraph{Core Predictors.}
Core predictors are metrics that are both consistently selected ($\geq$80\% frequency) \emph{and} have stable sign (all positive or all negative) across all mergers:

\begin{itemize}
    \item \textbf{ViT-B-16 (2 core predictors):}
    \begin{itemize}
        \item[$-$] \texttt{encoder\_gradient\_l2\_distance} (avg coef: $-$0.021)
        \item[$-$] \texttt{input\_gradient\_l2\_distance} (avg coef: $-$0.018)
    \end{itemize}
    \item \textbf{ViT-B-32 (3 core predictors):}
    \begin{itemize}
        \item[$-$] \texttt{encoder\_gradient\_l2\_distance} (avg coef: $-$0.026)
        \item[$+$] \texttt{activation\_magnitude\_ratio} (avg coef: $+$0.015)
        \item[$-$] \texttt{task\_vector\_l2\_distance} (avg coef: $-$0.008)
    \end{itemize}
\end{itemize}

The shared core predictor (\texttt{encoder\_gradient\_l2\_distance}) confirms that gradient-based metrics capture a fundamental, architecture-invariant signal for mergeability prediction.

\subsubsection{Sign Agreement Analysis}

Beyond metric selection, we analyze whether the \emph{direction} of effect (positive vs negative coefficient) is consistent across architectures. For each metric, we check if both models assign coefficients of the same sign across all mergers.

\textbf{8 metrics show 100\% sign agreement}, including:
\begin{itemize}
    \item \texttt{encoder\_gradient\_l2\_distance} (negative on both)
    \item \texttt{singular\_value\_overlap} (negative on both)
    \item \texttt{task\_vector\_l2\_distance} (negative on both)
    \item \texttt{right\_subspace\_overlap\_bottom\_k} (negative on both)
\end{itemize}

This sign consistency indicates that the \emph{interpretation} of metrics is stable: lower gradient distance, lower task vector distance, and lower singular value overlap all predict better merging on both architectures.

\subsubsection{Metric Category Ablation Analysis}

To quantify the importance of different metric categories, we perform ablation experiments where each category is entirely removed from the feature set and L1 LOTO optimization is re-run. The resulting drop in validation correlation measures category importance. Table~\ref{tab:category_ablation} summarizes the results.

\begin{table}[h]
\centering
\caption{Category ablation results: validation correlation drop ($\Delta$) when each metric category is removed. Negative values indicate the category is important for prediction.}
\label{tab:category_ablation}
\begin{tabular}{l|cc|cc|l}
\toprule
\textbf{Category} & \textbf{B-16 val $r$} & \textbf{B-16 $\Delta$} & \textbf{B-32 val $r$} & \textbf{B-32 $\Delta$} & \textbf{Importance} \\
\midrule
Baseline (all) & 0.544 & --- & 0.648 & --- & --- \\
\midrule
No GRAD\_BASED & 0.395 & $-$0.149 & 0.517 & $-$0.131 & \textbf{Critical} \\
No ACTIVATION & 0.546 & $+$0.002 & 0.589 & $-$0.059 & Mixed \\
No SUBSPACE & 0.499 & $-$0.045 & 0.634 & $-$0.014 & Moderate \\
No TASK\_VECTOR & 0.536 & $-$0.008 & 0.635 & $-$0.013 & Low \\
No EFF\_RANK & 0.564 & $+$0.020 & 0.660 & $+$0.012 & \textbf{Harmful} \\
\bottomrule
\end{tabular}
\end{table}

\paragraph{Main Conclusions Hold Across Architectures.}
The category ablation reveals two consistent findings:
\begin{enumerate}
    \item \textbf{GRAD\_BASED metrics are critical}: Removing gradient-based metrics causes the largest performance drop on both architectures ($-$0.149 on B-16, $-$0.131 on B-32). This confirms that \texttt{encoder\_gradient\_l2\_distance} and related metrics carry the primary predictive signal.
    \item \textbf{EFF\_RANK metrics introduce noise}: Removing effective rank metrics \emph{improves} performance on both architectures ($+$0.020 on B-16, $+$0.012 on B-32), indicating these metrics add noise rather than signal.
\end{enumerate}

The remaining categories (activation, subspace, task vector) show mixed results that vary by architecture, suggesting they provide secondary signals useful in specific contexts but are not universally critical.

\subsubsection{Discussion}

\paragraph{What Holds Across Architectures.}
\begin{enumerate}
    \item \textbf{Universal predictor}: \texttt{encoder\_gradient\_l2\_distance} is the most important metric on both architectures
    \item \textbf{Gradient dominance}: Gradient-based metrics receive the largest coefficients and are the most critical category (removing them causes $>$0.13 correlation drop)
    \item \textbf{Negative distance = better merging}: Distance metrics consistently have negative coefficients
    \item \textbf{Effective rank metrics are noise}: Removing them improves performance on both architectures
\end{enumerate}

\paragraph{Architecture-Specific Differences.}
\begin{enumerate}
    \item \textbf{Predictability}: ViT-B-32 shows higher validation correlations (0.648 vs 0.544), likely due to simpler feature spaces from larger patch sizes
    \item \textbf{Secondary predictors}: ViT-B-16 relies on \texttt{input\_gradient\_l2\_distance}, while ViT-B-32 uses \texttt{activation\_magnitude\_ratio} and \texttt{task\_vector\_l2\_distance}
    \item \textbf{Most predictable merger}: TSV is most predictable on ViT-B-16 (val $r$ = 0.627), while Weight Avg is most predictable on ViT-B-32 (val $r$ = 0.699)
\end{enumerate}

These differences are \emph{expected and desirable}---they demonstrate that the framework \emph{adapts} to architecture-specific characteristics rather than forcing a one-size-fits-all solution. The core predictive signal (gradient distance) is architecture-invariant, while secondary signals are tailored to each model.

\paragraph{Implications.}
The successful generalization to ViT-B-32 validates that our linear mergeability prediction framework:
\begin{itemize}
    \item Is not overfit to ViT-B-16's specific characteristics
    \item Identifies interpretable, architecture-invariant predictors (gradient metrics)
    \item Adapts secondary predictor selection to architecture-specific signals
    \item Can be applied to new architectures without methodology changes
\end{itemize}

This positions the framework as a practical tool for predicting merge compatibility across the ViT model family, with straightforward extension to other architectures and metrics.
